An \code{AddressSet} is a collection of \code{Process} and
\code{Dyninst::Address} tuples. \code{AddressSet} is used by the
ProcessSet\index{ProcessSet: : : :!} and ThreadSet\index{ThreadSet: : : :!}
classes for performing group operations on large numbers of processes. An
\code{AddressSet} might, for example, represent the location of a symbol across
numerous processes, or the location of a buffer in each process where data can
be written or read.  It satisfies the C++
\href{https://en.cppreference.com/cpp/named_req/ForwardIterator}{LegacyForwardIterator}
concept. It also offers the standard \code{begin/end} members for iteration.

\begin{apient}
AddressSet::ptr
AddressSet::const_ptr
\end{apient}

\apidesc{
The AddressSet::ptr and AddressSet::const\_ptr respectively represent a pointer
and a const pointer to an AddressSet object. Both pointer types are reference
counted and will cause the underlying AddressSet object to be cleaned when
there are no more references.
}

\begin{apient}
typedef std::pair<Dyninst::Address, Process::ptr> value_type
\end{apient}

\apidesc{
}

\begin{apient}
static AddressSet::ptr newAddressSet()
\end{apient}

\apidesc{
This function returns a new AddressSet that is empty.
}

\begin{apient}
static AddressSet::ptr newAddressSet(ProcessSet::const_ptr ps, Dyninst::Address addr)
\end{apient}

\apidesc{
This function returns a new AddressSet initialized with the elements from ps
paired with the Address addr.
}

\begin{apient}
static AddressSet::ptr newAddressSet(ProcessSet::const_ptr ps,
                                     std::string library_name,
                                     Dyninst::Offset off)
\end{apient}

\apidesc{
This function returns a new AddressSet initialized with the elements from ps.
The Address element for each process is calculated by looking up the load
address of library\_name in each Process and adding it to off.
}

\begin{apient}
iterator find(Dyninst::Address addr)
const_iterator find(Dyninst::Address addr) const
\end{apient}

\apidesc{
These functions return an iterator that points to the first element in the
AddressSet with an address of addr. They return end() if no element matches
addr.
}

\begin{apient}
iterator find(Dyninst::Address addr, Process::const_ptr proc)
const_iterator find(Dyninst::Address addr, Process::const_ptr proc) const
\end{apient}

\apidesc{
These functions return an iterator that points to any element that has a process
and address of proc and addr. It returns end() if no element matches.
}

\begin{apient}
size_t count(Dyninst::Address addr) const
\end{apient}

\apidesc{
This function returns the number of elements with address addr.
}

\begin{apient}
size_t size() const
\end{apient}

\apidesc{
This function returns the number of elements in the AddressSet.
}

\begin{apient}
bool empty() const
\end{apient}

\apidesc{
This function returns true if the AddressSet has zero elements and false
otherwise.
}

\begin{apient}
std::pair<iterator, bool> insert(Dyninst::Address addr, Process::const_ptr proc)
\end{apient}

\apidesc{
This function inserts a new element into the AddressSet with addr and proc as
its values. If another element with those values already exists, then no new
element will be inserted. It returns an iterator that points to the new or
existing element and a boolean value that is true if a new element was inserted
and false otherwise.
}

\begin{apient}
size_t insert(Dyninst::Address addr, ProcessSet::const_ptr ps)
\end{apient}

\apidesc{
For every element in ps, this function inserts it and addr into the AddressSet.
It returns the number of new elements created.
}

\begin{apient}
void erase(iterator pos)
\end{apient}

\apidesc{
This function removes the element pointed to by pos from the AddressSet.
}

\begin{apient}
size_t erase(Process::const_ptr proc)
\end{apient}

\apidesc{
This function removes every element with a process of proc from the AddressSet.
It returns the number of elements removed.
}

\begin{apient}
size_t erase(Dyninst::Address addr,
\end{apient}

\apidesc{
Process::const\_ptr proc)
This function removes any element that has and address and process of addr and
proc from the AddressSet. It returns the number of elements removed.
}

\begin{apient}
void clear()
\end{apient}

\apidesc{
This function erases all elements from the AddressSet leaving an AddressSet of
size zero.
}

\begin{apient}
iterator lower_bound(Dyninst::Address addr)
\end{apient}

\apidesc{
This function returns an iterator pointing to the first element in the
AddressSet that has an address greater than or equal to addr.
}

\begin{apient}
iterator upper_bound(Dyninst::Address addr)
\end{apient}

\apidesc{
This function returns an iterator pointing to the first element in the
AddressSet that has an address greater than addr.
}

\begin{apient}
std::pair<iterator, iterator> equal_range(Address addr) const
\end{apient}

\apidesc{
This function returns a pair of iterators. The first iterator has the same
value as the return of lower\_bound(addr) and the second iterator has the same
value as the return of upper\_bound(addr).
}

\begin{apient}
AddressSet::ptr set_union(AddressSet::const_ptr aset)
\end{apient}

\apidesc{
This function returns a new AddressSet whose elements are the set union of this
AddressSet and aset.
}

\begin{apient}
AddressSet::ptr set_intersection(AddressSet::const_ptr aset)
\end{apient}

\apidesc{
This function returns a new AddressSet whose elements are the set intersection
of this AddressSet and aset.
}

\begin{apient}
AddressSet::ptr set_difference(AddressSet::const_ptr aset)
\end{apient}

\apidesc{
This function returns a new AddressSet whose elements are the set difference of
this AddressSet minus aset.
}
